% =================================================
% D 册补充专题：孕期性生活 / 更年期性生活
% 现代循证撰写，风格对齐本册既有章节
% =================================================

\chapter{孕期性生活专题}

\section{先说结论}
在\textbf{无并发症}的妊娠中，性生活对母婴一般是\textbf{安全的}，无需整个孕期"禁欲"。
胎盘和充满羊水的羊膜囊会保护胎儿，正常性接触不会压到、碰伤宝宝，也不会引起早产。
真正需要调整或暂停的，是\textbf{有特定医学指征}的妊娠——这类情形由产科医生个案判断，而非一刀切。

\section{各孕期的关注点}
\begin{tabularx}{\textwidth}{l X}
\toprule
\textbf{阶段} & \textbf{关注点} \\
\midrule
孕早期（0--12 周） & 多数人因孕吐、疲劳而性欲下降，属正常；如有流产史或先兆流产（出血、腹痛），遵医嘱暂停。 \\
孕中期（13--27 周） & 常为最舒适的一段：孕吐减轻、精力恢复，性欲可能回升；姿势以不压迫腹部为宜。 \\
孕晚期（28 周--临产） & 腹部增大，宜侧卧或从后进入，避免仰卧位压迫下腔静脉；临产前数日如胎盘位置低、已破水或有出血，暂停。 \\
\bottomrule
\end{tabularx}

\section{这些情况请暂停并咨询产科医生}
\begin{itemize}
  \item 前置胎盘、不明原因阴道出血；
  \item 宫颈机能不全、已放置宫颈环扎（cerclage）；
  \item 早破水（胎膜早破）或已破水；
  \item 先兆早产、持续规律宫缩；
  \item 医生已明确建议"卧床 / 暂停性生活"的其他情形。
\end{itemize}

\section{体位与卫生的实用建议}
\begin{itemize}
  \item \textbf{姿势}：孕晚期首选侧卧位、女上位或背后位，避免长时间仰卧压迫腹部与下腔静脉；动作放慢、减少腹部受压。
  \item \textbf{润滑}：孕期阴道分泌增多、润滑通常更好；若偏干可用\textbf{水溶性润滑剂}。
  \item \textbf{安全套}：孕中后期性生活一般不会再次受孕，但为减少泌尿生殖道感染、并避免精液前列腺素可能诱发宫缩，部分医生建议使用。
  \item \textbf{射后}：射入或射在会阴都属正常，无需刻意冲洗阴道内部。
  \item \textbf{卫生}：性生活前后洗手、清洁外阴即可，不必灌洗阴道。
\end{itemize}

\section{常见疑问}
\begin{itemize}
  \item \textbf{会让宝宝"缺氧 / 挤压"吗？} 不会。羊水与子宫肌层提供缓冲，正常强度下无碍。
  \item \textbf{会早产吗？} 无并发症者无证据表明正常性生活致早产；有早产风险者遵医嘱。
  \item \textbf{高潮会伤胎吗？} 宫缩是高潮的短暂生理反应，通常几分钟内消退，无并发症者安全。
  \item \textbf{丈夫可以口交吗？} 一般可以，但应避免向阴道内吹气（极罕见的空气栓塞风险），并避免经口传播感染。
\end{itemize}

\section{一页要点}
无并发症的妊娠，性生活对母婴安全；有出血、先兆流产、宫颈环扎、早破水、前置胎盘或早产风险时，\textbf{听产科医生的个案建议}。
孕晚期优先侧卧 / 背后位、避免仰卧压腹，注意润滑与卫生。有疑问先问医生，不要凭"老说法"一刀切禁欲。


\chapter{更年期性生活专题}

\section{先正名：更年期不等于"没性欲了""}
围绝经期（perimenopause，多在 45--55 岁）到绝经后，雌激素下降会带来身体变化，
但\textbf{性欲和能力并不会随绝经自动消失}。很多女性报告高潮强度反而增强、不再担心怀孕、更放松投入。
真正的困扰多来自\textbf{具体可处理的症状}（干涩、疼痛、情绪波动），而非"老了就不行了"的观念。

\section{身体上会发生什么，以及对策}
\begin{tabularx}{\textwidth}{l X}
\toprule
\textbf{变化} & \textbf{应对} \\
\midrule
阴道干涩、弹性下降（genitourinary syndrome of menopause, GSM） & 充分前戏 + \textbf{水溶性润滑剂}；反复干涩疼痛可请医生评估局部雌激素。 \\
唤起变慢、需要更直接的刺激 & 延长前戏，增加阴蒂刺激，不必以插入为唯一目标。 \\
潮热 / 夜汗、情绪波动 & 规律作息与运动；影响睡眠者可咨询医生（含激素治疗 HT 的利弊）。 \\
皮肤与黏膜更敏感 & 动作放缓，避免过紧过快的摩擦。 \\
\bottomrule
\end{tabularx}

\section{关于激素治疗（HT）的一句话}
局部雌激素软膏 / 栓剂对\textbf{阴道干涩与性交痛}证据充分、风险低，是围绝经期的常用一线。
是否加用全身激素治疗需\textbf{个体化}权衡（年龄、症状、心血管与乳腺风险），务必与医生共同决定，不要自行长期用药。

\section{关系与心理：常被忽略的一半}
\begin{itemize}
  \item 更年期常伴随"孩子离家（空巢）"，部分夫妻反而进入\textbf{更自由、更专注彼此}的阶段。
  \item 性欲差异在围绝经期可能被放大（女性因激素下降、伴侣因其他原因），核心仍是\textbf{沟通与共同探索}，而非"谁该迁就谁"。
  \item 把关注点从"插入与高潮"扩展到\textbf{爱抚、亲密、情感联结}，往往比追求次数更能维持满意。
\end{itemize}

\section{一页要点}
更年期不是性生活的终点。干涩、疼痛等具体症状\textbf{大多可处理}：润滑剂 + 充分前戏是基础，持续困扰请看医生评估局部雌激素。
心态上，把性从"任务"还原为"亲密"，多数女性能在新阶段找到适合自己的方式。有疑问，找妇科 / 更年期门诊，而不是默认忍到"算了"。


\chapter{性侵后的妇科处理与创伤细节（专科补充）}

\section{呼应 A 册，这里讲"女性身体"的专科部分}
A 册《性侵犯与强奸》已覆盖认知、急救取证与通用健康后果。本专册补充\textbf{女性特有的}妇科检查、损伤处理与后续。
核心结论先说：\textbf{绝大多数伤情可处理，感染可预防，妊娠可应对，心理可恢复}——关键是\textbf{尽快、按自己节奏}去妇科/急诊，别硬扛。

\section{妇科检查会做什么（心里有数，不慌张）}
\begin{itemize}
  \item \textbf{视诊 + 触诊}：外阴、阴道口、会阴有无撕裂/血肿/擦伤；必要时肛诊。动作尽量轻柔，可要求女医护陪同。
  \item \textbf{取样检测（rape kit）}：棉拭子取阴道/肛道/口腔分泌物，查精液、皮屑、纤维及 STI（HIV、梅毒、淋病、衣原体、HSV、HPV 等）。
  \item \textbf{妊娠评估}：判断是否需紧急避孕（见下）。
  \item \textbf{伤情拍照/记录}：淤青、抓痕、撕裂留档，供医疗与（若需要）法律使用。
\end{itemize}
\textbf{提示}：未就医前尽量避免冲洗阴道内部、擦洗伤口，以免冲掉物证；若已清洗，如实告知医生即可。

\section{女性特有的处理要点}
\begin{tabularx}{\textwidth}{l X}
\toprule
\textbf{问题} & \textbf{处理} \\
\midrule
阴道/外阴撕裂、血肿 & 急诊清创缝合；血肿大时压迫 + 必要时手术。 \\
泌尿道感染 / 损伤 & 血尿、尿频尿痛查尿常规与培养，按结果抗感染。 \\
STI 暴露 & HIV PEP（\textbf{72 小时内}最有效）、梅毒/淋病/衣原体预防性用药、后续 2--4 周/3 月/6 月复测。 \\
意外妊娠 & 紧急避孕：口服紧急避孕药（72h 内）或\textbf{铜宫内节育器（5 天内，成功率更高}）；后续选常规避孕。 \\
HPV / 宫颈癌筛查 & 事后纳入常规随访；持续感染者按医嘱管理。 \\
慢性盆腔痛 / 性交痛（dyspareunia） & 多与盆底肌紧张 + 创伤相关，可做盆底康复、理疗 + 心理干预，可治。 \\
\bottomrule
\end{tabularx}

\section{对"事后羞耻/不敢复查"的一句提醒}
\begin{itemize}
  \item 复查、取环、做 HPV 筛查，和任何一次妇科常规一样\textbf{平常}，不用向医生解释前情。
  \item 若出现\textbf{性交痛、经量异常、反复尿路感染、盆底痛}，及时就诊——这些\textbf{不是"忍忍就好"的常态}。
  \item 心理与身体可以分开处理：先解决感染/妊娠/伤情，再处理情绪；也可以同步做。
  \item 你有权决定"做不做、什么时候做、跟谁说"，\textbf{节奏由你定}，而不是被旁人催着"该放下了"。
\end{itemize}

\section{一页要点}
女性性侵后的专科重点：妇科检查（rape kit 取样）+ 损伤处理 + STI/HPV 暴露预防 + 紧急避孕 + 盆底康复。
72 小时内的 PEP 窗口、5 天内的铜 IUD 紧急避孕，是最需要"赶早"的两件事。
持续性交痛/盆腔痛/反复尿路感染要复查，别当成常态硬扛。节奏由你自己定。


\chapter{附录：老书章要点速查（D 册收录的旧书整理章）}
以下各章为\textbf{1990s 旧书}整理，核心信息仍有效，但药物/数据请以下篇 D 册"现代循证章"为准：

\begin{itemize}
  \item \textbf{女性生殖系统}：解剖、月经周期、排卵机制；\textbf{图式偏旧}，建议结合 D 册上篇"女性生理"章对照。
  \item \textbf{女性在性交中的意外（旧书整理）}：性交痛、阴道痉挛、经期/孕期性生活的注意事项；\textbf{核心有效}，孕期部分与 D 册《孕期性生活专题》重叠，以后者为准。
  \item \textbf{妇女保健及病症治疗}：妇科常见症状（白带异常、经期不适、更年期）的识别与就医；\textbf{部分药物已过时}，现代指南（ACOG/WHO）已更新。
  \item \textbf{女性不孕与治疗（旧书整理）}：不孕原因排查（排卵/输卵管/激素）、治疗路径；\textbf{核心有效}，辅助生殖（IVF/ICSI）技术进展快，详见 D 册上篇。
  \item \textbf{人工流产、引产知识}：人流/引产的方法、术后恢复、注意事项；\textbf{技术仍主流}，但药物流产比例已大幅提高，详见 D 册上篇。
  \item \textbf{女性绝育术}：输卵管结扎术原理、术后注意、后悔率；\textbf{技术仍主流}，"不可逆"程度比旧书描述更可控（复通/取卵仍有可能）。
\end{itemize}

\textbf{阅读建议}：旧书章的\textbf{生理机制与就医路径}仍有效，\textbf{具体药物/技术数据}以下篇循证章为准。


\chapter{盆底康复与孕产期情趣用品}

\begin{center}
\myhei\large 盆底是"被低估的开关"——练对了，疼痛、控制、敏感度都会跟着变好
\end{center}

\section{盆底肌为什么重要}
盆底肌是一组托住膀胱、子宫、直肠的"吊床"肌肉，功能包括：
\begin{itemize}[leftmargin=2em,itemsep=0.3em]
  \item \textbf{承托}：托住盆腔脏器，预防脱垂。
  \item \textbf{控尿}：协助尿道关闭，减少漏尿（尤其咳嗽、跳跃时）。
  \item \textbf{性功能}：高潮的节律性收缩、敏感度、控制感都与盆底状态相关（见 B 册《享受高潮》章"呼吸与盆底"一节）。
  \item \textbf{分娩/产后}：分娩会拉伸盆底，产后恢复决定长期功能。
\end{itemize}
\textbf{一句话}：盆底不是"女人的事"，男女都有，且都受益。

\section{盆底训练器类型}
\begin{tabularx}{\textwidth}{X X X}
\toprule
\textbf{类型} & \textbf{原理} & \textbf{适合} \\
\midrule
Kegel 器械（重锤/锥形塞） & 负重被动训练，让肌肉"顶住"重量 & \textbf{入门}、产后恢复期（需医生确认时机） \\
生物反馈（biofeedback） & 实时显示肌肉收缩/放松，学"找对肌肉" & \textbf{不知道"怎么练"}、练了没效果的人 \\
电刺激（electrostimulation） & 电流刺激肌肉被动收缩 & \textbf{肌力弱、自主收缩困难}者；需专业评估 \\
智能 App 器械 & 结合传感器 + App 指导 & \textbf{喜欢"打卡"}、需要反馈的人 \\
\bottomrule
\end{tabularx}
\par\smallskip
\textbf{关键}：\textbf{先"找对肌肉"再谈器械}——很多人用器械却没练到盆底（练的是臀肌/腹肌），\textbf{生物反馈是最好的"老师"}。

\section{孕晚期：适用与禁忌}
\begin{itemize}[leftmargin=2em,itemsep=0.3em]
  \item \textbf{盆底训练器}：孕晚期\textbf{以"被动放松"为主}，\textbf{别用重锤}（会增加腹压）；\textbf{电刺激需医生评估}。
  \item \textbf{情趣玩具}：\textbf{避免"深插入 + 强振动"}（可能刺激宫缩）；\textbf{侧卧位}最安全；\textbf{干涩用润滑}（孕期激素变化常致干涩）。
  \item \textbf{禁忌情况}：前置/低位胎盘、破水、不明原因出血、多次流产、宫颈机能不全——\textbf{先问医生}。
  \item \textbf{原则}：\textbf{"不舒服就停"}，\textbf{疼就就医}（不是"忍一忍"）。
\end{itemize}

\section{产后恢复期}
\begin{itemize}[leftmargin=2em,itemsep=0.3em]
  \item \textbf{时机}：\textbf{顺产 6 周、剖宫产 6--8 周}（视伤口恢复），\textbf{先复查}再恢复盆底训练和性生活。
  \item \textbf{顺序}：\textbf{先"放松"再"收紧"}（产后盆底常"紧张"而非"松弛"，直接练收紧可能加重症状）；\textbf{生物反馈是首选}。
  \item \textbf{情趣用品}：\textbf{干涩是常态}（哺乳期激素低）——\textbf{润滑剂 + 保湿剂}组合；\textbf{别急着"恢复从前"}，身体在适应新状态。
  \item \textbf{漏尿/脱垂}：\textbf{盆底康复科}评估，必要时\textbf{电刺激 + 生物反馈 + 行为训练}联合。
\end{itemize}

\section{就医路径（找谁）}
\begin{enumerate}[leftmargin=2em,itemsep=0.3em]
  \item \textbf{盆底康复科/康复医学科}：系统评估（肌力、张力、协调）+ 制定训练方案。
  \item \textbf{妇科}：脱垂、漏尿、产后复查。
  \item \textbf{泌尿科（女性}）：压力性尿失禁的进一步评估。
  \item \textbf{怎么判断"需要专业帮助"}：\textbf{自我训练 4--6 周无改善}、\textbf{漏尿影响生活}、\textbf{产后 6 周仍疼痛/无力}——\textbf{别硬练}，找专业评估。
\end{enumerate}

\section{一页要点}
盆底 = \textbf{承托 + 控尿 + 性功能}的"三合一"开关，男女都受益。
\textbf{训练器}：重锤（入门）/ 生物反馈（找对肌肉）/ 电刺激（肌力弱）/ 智能 App（打卡）——\textbf{先"找对肌肉"再上器械}。
\textbf{孕晚期}：以"放松"为主、别用重锤、侧卧位、干涩用润滑；\textbf{前置胎盘/破水/出血先问医生}。
\textbf{产后}：顺产 6 周/剖宫产 6--8 周再恢复；\textbf{先放松再收紧}；干涩用润滑+保湿；漏尿/脱垂找盆底康复科。
\textbf{就医}：盆底康复科/妇科/泌尿科——\textbf{自我训练 4--6 周无改善就找专业评估}，别硬练。
